\subsection{2.6 Das universelle Stabilitätsfunktional}

Nachdem sich gezeigt hatte, dass der rohe OPP keine intrinsische vierdimensionale Raumzeit besitzt, stellte sich zwangsläufig eine neue Frage:

\begin{center}
\emph{Existieren innerhalb des hochdimensionalen OPP bevorzugte stabile Zustände?}
\end{center}

Mit anderen Worten:

\[
\boxed{
\text{Kann Dynamik allein bestimmte Strukturen bevorzugen?}
}
\]

Diese Überlegung führte zur Einführung eines universellen Stabilitätsfunktionals, das unabhängig von einer konkreten physikalischen Theorie formuliert werden sollte.

%%%%%%%%%%%%%%%%%%%%%%%%%%%%%%%%%%%%%%%%%%%%%%%%%%%%%%%%%%%%%%

\section*{Grundidee}

Jedes komplexe System besitzt typischerweise konkurrierende Anforderungen.

Einerseits möchte das System möglichst viele Freiheitsgrade besitzen.

Andererseits erzeugen zu viele Freiheitsgrade Instabilität.

Auch im OPP wurde angenommen, dass zwei Größen miteinander konkurrieren:

\begin{enumerate}

\item maximale Selbstinformation

\item minimale Frustration

\end{enumerate}

Die gesuchte Dynamik sollte daher nicht möglichst viele Freiheitsgrade erzeugen,

sondern einen optimalen Kompromiss finden.

%%%%%%%%%%%%%%%%%%%%%%%%%%%%%%%%%%%%%%%%%%%%%%%%%%%%%%%%%%%%%%

\section*{Selbstinformation}

Als erste Größe wurde die gesamte Selbstinformation des Systems definiert.

Sie wurde durch die Spur der projektiven Gram-Matrix beschrieben:

\[
\Phi
=
\operatorname{Tr}(K).
\]

Dabei bezeichnet

\[
K
\]

die aus den Relationen konstruierte Gram-Matrix.

Große Werte von

\[
\Phi
\]

entsprechen einer hohen Gesamtinformation des Systems.

%%%%%%%%%%%%%%%%%%%%%%%%%%%%%%%%%%%%%%%%%%%%%%%%%%%%%%%%%%%%%%

\section*{Frustration}

Als zweite Größe wurde die sogenannte Frustration eingeführt.

Die Idee stammt ursprünglich aus der Physik komplexer Systeme und der Spin-Gläser.

Ein ideales System besitzt nur wenige dominante Eigenmoden.

Viele kleine Eigenwerte bedeuten dagegen,

dass Information auf viele Freiheitsgrade verteilt wird.

Die Frustration wurde daher definiert als

\[
F
=
\sum_{i>r}
\lambda_i^2,
\]

wobei

\[
r
\]

den gewünschten Zielrang bezeichnet.

Je größer die Beiträge oberhalb dieses Zielranges,

desto stärker ist die Frustration.

%%%%%%%%%%%%%%%%%%%%%%%%%%%%%%%%%%%%%%%%%%%%%%%%%%%%%%%%%%%%%%

\section*{Das Stabilitätsfunktional}

Aus beiden Größen wurde schließlich das universelle Funktional konstruiert:

\[
\boxed{
\mathcal{S}
=
\Phi
-
\lambda F.
}
\]

Dabei beschreibt

\[
\lambda
\]

die relative Gewichtung der Frustration.

Das Ziel der Optimierung lautet

\[
\mathcal{S}
\rightarrow
\max.
\]

%%%%%%%%%%%%%%%%%%%%%%%%%%%%%%%%%%%%%%%%%%%%%%%%%%%%%%%%%%%%%%

\section*{Physikalische Interpretation}

Dieses Funktional besitzt eine bemerkenswert einfache Bedeutung.

Der erste Term

\[
\Phi
\]

belohnt möglichst viel Information.

Der zweite Term

\[
F
\]

bestraft unnötig komplexe Eigenwertspektren.

Das System versucht also,

möglichst viel Information

mit möglichst wenigen dominanten Freiheitsgraden

zu speichern.

%%%%%%%%%%%%%%%%%%%%%%%%%%%%%%%%%%%%%%%%%%%%%%%%%%%%%%%%%%%%%%

\section*{Numerische Optimierung}

Zur Maximierung wurde ein Standardverfahren verwendet.

Zum Einsatz kam insbesondere

\[
\text{L-BFGS-B}
\]

als numerischer Optimierer.

Ausgehend von zufälligen Punktkonfigurationen wurden die Attraktoren solange verschoben,

bis

\[
\mathcal{S}
\]

ein lokales Maximum erreichte.

%%%%%%%%%%%%%%%%%%%%%%%%%%%%%%%%%%%%%%%%%%%%%%%%%%%%%%%%%%%%%%

\section*{Das überraschende Ergebnis}

Bereits die ersten Optimierungen lieferten erstaunliche Resultate.

Während der rohe OPP typischerweise

\[
d_{\mathrm{eff}}
\approx
30\ldots45
\]

besaß,

ergaben optimierte Konfigurationen häufig

\[
d_{\mathrm{eff}}
\approx
4\ldots5.
\]

Dies war die erste numerische Beobachtung,

die überhaupt in die Nähe einer beobachtbaren Raumzeit führte.

%%%%%%%%%%%%%%%%%%%%%%%%%%%%%%%%%%%%%%%%%%%%%%%%%%%%%%%%%%%%%%

\section*{Der scheinbare Durchbruch}

Die ersten Resultate wirkten zunächst äußerst spektakulär.

Insbesondere traten mehrfach Werte auf wie

\[
d_{\mathrm{eff}}
=
3.75,
\]

oder

\[
d_{\mathrm{eff}}
=
4.1.
\]

Es entstand kurzfristig der Eindruck,

dass das Stabilitätsfunktional tatsächlich spontan eine Raumzeit erzeugen könnte.

%%%%%%%%%%%%%%%%%%%%%%%%%%%%%%%%%%%%%%%%%%%%%%%%%%%%%%%%%%%%%%

\section*{Die kritische Analyse}

Im weiteren Verlauf wurde das Verfahren jedoch sorgfältig hinterfragt.

Dabei zeigte sich ein entscheidendes Problem.

Die Frustration war definiert als

\[
F
=
\sum_{i>r}
\lambda_i^2.
\]

Der Zielrang

\[
r
\]

war also bereits Bestandteil des Funktionals.

Damit wurde implizit bereits vorgegeben,

welche Anzahl dominanter Eigenwerte bevorzugt wird.

Dies bedeutet:

\[
\boxed{
\text{Das Verfahren war nicht vollständig rangneutral.}
}
\]

%%%%%%%%%%%%%%%%%%%%%%%%%%%%%%%%%%%%%%%%%%%%%%%%%%%%%%%%%%%%%%

\section*{Warum das trotzdem wichtig war}

Diese Erkenntnis bedeutete keineswegs,

dass die Optimierung wertlos gewesen wäre.

Ganz im Gegenteil.

Sie zeigte erstmals,

dass sich Eigenwertspektren des OPP gezielt beeinflussen lassen.

Insbesondere wurde deutlich,

dass der OPP keineswegs starr ist,

sondern auf Variationsprinzipien reagiert.

%%%%%%%%%%%%%%%%%%%%%%%%%%%%%%%%%%%%%%%%%%%%%%%%%%%%%%%%%%%%%%

\section*{Die wissenschaftliche Konsequenz}

Anstatt das Resultat vorschnell als Bestätigung der PST zu interpretieren,

wurde ein neuer Forschungsweg eingeschlagen.

Die zentrale Frage lautete nun:

\[
\boxed{
\text{Welche Resultate bleiben erhalten, wenn keinerlei Zielrang mehr vorgegeben wird?}
}
\]

Diese Fragestellung führte unmittelbar zur nächsten Phase der Untersuchungen.

%%%%%%%%%%%%%%%%%%%%%%%%%%%%%%%%%%%%%%%%%%%%%%%%%%%%%%%%%%%%%%

\section*{Neue Teststrategie}

Von diesem Zeitpunkt an wurden sämtliche Optimierungsverfahren zusätzlich durch rangneutrale Methoden überprüft.

Dazu gehörten insbesondere

\begin{itemize}

\item verschiedene Projektionsverfahren,

\item Kernel-Vergleiche,

\item Diffusionsoperatoren,

\item Random-Projektionen,

\item universelle Skalierungstests.

\end{itemize}

Nur Resultate,

die auch ohne eingebaute Ranginformation auftraten,

wurden anschließend als mögliche physikalische Hinweise betrachtet.

%%%%%%%%%%%%%%%%%%%%%%%%%%%%%%%%%%%%%%%%%%%%%%%%%%%%%%%%%%%%%%

\section*{Bedeutung für die PST}

Rückblickend markiert das Stabilitätsfunktional einen wichtigen Wendepunkt.

Es zeigte,

dass

\begin{enumerate}

\item Variationsprinzipien grundsätzlich geeignet sind,

Eigenwertspektren zu strukturieren,

\item aber gleichzeitig,

dass jede implizite Vorannahme sorgfältig kontrolliert werden muss.

\end{enumerate}

Gerade diese Selbstkritik machte das gesamte numerische Forschungsprogramm wesentlich robuster.

%%%%%%%%%%%%%%%%%%%%%%%%%%%%%%%%%%%%%%%%%%%%%%%%%%%%%%%%%%%%%%

\section*{Zusammenfassung}

Das universelle Stabilitätsfunktional war der erste Versuch,

aus rein informationstheoretischen Prinzipien stabile Strukturen innerhalb des OPP zu erzeugen.

Obwohl sich später zeigte,

dass der Zielrang bereits teilweise eingebaut war,

lieferte das Verfahren dennoch mehrere wichtige Erkenntnisse:

\begin{itemize}

\item Der OPP reagiert auf Variationsprinzipien.

\item Dominante Eigenmoden können gezielt verstärkt werden.

\item Die Optimierung allein erklärt jedoch noch keine emergente Raumzeit.

\item Eine echte Erklärung muss ohne implizite Rangvorgaben funktionieren.

\end{itemize}

Diese Einsicht führte unmittelbar zur entscheidenden nächsten Entwicklungsstufe der PST:

der systematischen Untersuchung verschiedener mathematischer Projektionsoperatoren.

